\documentclass[12pt, a4paper]{article}

% ── Pacotes ─────────────────────────────────────────────────────
\usepackage[utf8]{inputenc}
\usepackage[T1]{fontenc}
\usepackage[brazil]{babel}
\usepackage[top=2.5cm, bottom=2.5cm, left=2.5cm, right=2.5cm]{geometry}
\usepackage{setspace}
\usepackage{indentfirst}
\usepackage{titlesec}
\usepackage{enumitem}
\usepackage{hyperref}
\usepackage{parskip}
\usepackage{times}

\onehalfspacing

% Remove indentação de parágrafos (estilo template Word)
\setlength{\parindent}{0pt}
\setlength{\parskip}{6pt}

% Bullet personalizado ▪ para listas
\setlist[itemize,1]{label=\textbullet, leftmargin=1.5em}

% Formatação de seções
\titleformat{\section}{\normalfont\bfseries\large}{\thesection.}{0.5em}{}
\titleformat{\subsection}{\normalfont\bfseries}{\thesubsection.}{0.5em}{}
\titlespacing{\section}{0pt}{12pt}{4pt}
\titlespacing{\subsection}{0pt}{8pt}{2pt}

\hypersetup{colorlinks=true, linkcolor=black, citecolor=black, urlcolor=blue}

% ================================================================
\begin{document}
% ================================================================

% ── Cabeçalho do template ────────────────────────────────────────
\begin{center}
{\large\bfseries AA01: Proposta de Pesquisa \& Levantamento Bibliográfico Inicial}\\[0.3em]
\textit{Disciplina: Metodologia de Pesquisa}\\
\textit{Profa: Glívia Angélica Rodrigues Barbosa}\\
\textit{Departamento de Computação (DECOM)}\\
\textit{Centro Federal de Educação Tecnológica de Minas Gerais (CEFET-MG)}
\end{center}

\vspace{0.8em}
\underline{\textit{\textbf{Nome:} Morais, Ester}}

\vspace{1em}

% ================================================================
\section{Proposta}
% ================================================================

\begin{itemize}

  \item \textbf{Área e Tema/Demanda a ser explorada:}
  Área de \textbf{Segurança de Software}, com foco em Teste Dinâmico de Segurança de Aplicações (\textit{Dynamic Application Security Testing} -- DAST).
  O tema é a \textbf{detecção automatizada de vulnerabilidades em \textit{Single Page Applications} (SPAs) por meio de agentes de Inteligência Artificial baseados em Modelos de Linguagem de Grande Porte (LLMs)}.

  \item \textbf{Questão a ser explorada:}
  Em que medida um agente baseado em LLM, equipado com automação de navegador e orientado pelo padrão de raciocínio ReAct, é capaz de detectar vulnerabilidades em SPAs com maior eficácia (Precisão, Revocação e F1-\textit{Score}) do que ferramentas DAST convencionais?

  \item \textbf{Motivação:}
  Ferramentas DAST tradicionais, como o OWASP ZAP, operam por varredura de URLs e injeção automatizada de \textit{payloads}, pressupondo que a superfície de ataque é estática e mapeável por hiperlinks.
  Entretanto, SPAs construídas com \textit{frameworks} reativos (Angular, React, Vue) renderizam rotas dinamicamente via JavaScript, tornando grande parte dos estados da aplicação invisíveis ao \textit{spider} convencional.
  Adicionalmente, certas categorias de vulnerabilidade -- como falhas de lógica de negócio (ex.: quantidade negativa em um carrinho de compras), controle de acesso quebrado por identificadores previsíveis (IDOR) e ataques a tokens JWT -- exigem compreensão semântica do domínio, que está além das capacidades de \textit{scanners} baseados em regras.
  LLMs emergem como candidatos naturais para suprir essa lacuna, pois demonstraram capacidade de raciocinar sobre código, protocolos HTTP e fluxos de negócio.

  \item \textbf{Objetivo:}
  Projetar, implementar e avaliar um \textit{framework} de DAST agêntico -- denominado \textbf{Agentic DAST} -- que utiliza LLMs como motor cognitivo para navegar autonomamente os estados internos de uma SPA, raciocinar sobre fluxos de negócio e executar casos de abuso predefinidos.
  O \textit{framework} deve ser capaz de detectar múltiplas categorias de vulnerabilidade presentes no OWASP Top 10: lógica de negócio, controle de acesso quebrado, \textit{Cross-Site Scripting} (XSS), ataques a JWT, exposição de dados sensíveis e injeção de SQL.

  \item \textbf{Justificativa/Relevância:}
  A relevância da pesquisa é dupla.
  Do ponto de vista \textbf{científico}, o trabalho contribui para o campo emergente de \textit{Agentic AI for Security} ao propor uma arquitetura concreta que combina o padrão ReAct~\cite{yao2023react} com automação de navegador via \textit{Chrome DevTools Protocol} (CDP) e uma biblioteca de Casos de Abuso formalmente especificados -- uma combinação inédita na literatura revisada.
  Do ponto de vista \textbf{prático}, organizações que desenvolvem SPAs carecem de ferramentas automatizadas capazes de detectar falhas semânticas antes da implantação em produção; o \textit{framework} proposto preenche parte dessa lacuna de forma reproduzível e extensível, utilizando o OWASP Juice Shop como \textit{testbed} amplamente reconhecido pela comunidade.

  \item \textbf{Estilo da Pesquisa:}
  \begin{itemize}
    \item[(a)] \textbf{Definição de Wazlawick (2014) para Computação:} Pesquisa \textit{tecnológica} orientada ao desenvolvimento de artefato. O produto central da investigação é o \textit{framework} Agentic DAST; sua avaliação empírica valida as afirmações de desempenho~\cite{wazlawick2014}.
    \item[(b)] \textbf{Natureza:} Aplicada -- o problema origina-se de uma demanda real da indústria de segurança de \textit{software}.
    \item[(c)] \textbf{Abordagem:} Quantitativa -- os resultados são medidos por Precisão, Revocação e F1-\textit{Score}, calculados a partir de uma matriz \textit{Ground Truth} derivada do catálogo oficial de desafios do OWASP Juice Shop.
    \item[(d)] \textbf{Objetivo do Estilo:} Experimental -- compara o \textit{framework} proposto (Cenário~C) contra dois \textit{baselines}: (A)~OWASP ZAP com \textit{spider} tradicional; e (B)~ZAP guiado por script E2E de navegação humana.
    \item[(e)] \textbf{Procedimento:} Estudo experimental controlado com um único \textit{testbed} (OWASP Juice Shop em contêiner Docker), três rodadas de execução por cenário ($n = 3$) e análise estatística não-paramétrica (teste de Wilcoxon com $\alpha = 0{,}05$) para comparação de médias de F1-\textit{Score} entre os cenários.
  \end{itemize}

\end{itemize}

% ================================================================
\section{Trabalhos Relacionados}
% ================================================================

\subsection{Trabalho: PentestGPT: An LLM-Empowered Automatic Penetration Testing Framework}

\begin{itemize}

  \item \textbf{Referência:}
  DENG, Gelei; LIU, Yi; MAYORAL-VILCHES, Victor; LIU, Peng; LI, Yuekang; XU, Yuan; ZHANG, Tianwei; LIU, Yang; PINZGER, Martin; RASS, Stefan.
  PentestGPT: An LLM-Empowered Automatic Penetration Testing Framework.
  In: \textit{Proceedings of the USENIX Security Symposium}, 2024.
  Disponível em: \url{https://arxiv.org/abs/2308.06782}.

  \item \textbf{Resumo:}

  \textit{Motivação e Objetivo:} Testes de penetração manuais são processos custosos, demorados e altamente dependentes da expertise do profissional, o que limita sua adoção em organizações com menor capacidade técnica ou orçamentária.
  O trabalho propõe o PentestGPT, um \textit{framework} que emprega LLMs como motor de raciocínio para conduzir testes de penetração de forma autônoma ou semiautônoma, reduzindo a barreira de entrada para esse tipo de avaliação de segurança.

  \textit{Metodologia:} O \textit{framework} divide o processo de \textit{pentesting} em três módulos especializados e interligados: (1)~\textit{Reasoning Module}, responsável por traçar estratégias de alto nível com base no contexto acumulado; (2)~\textit{Generation Module}, que traduz estratégias em comandos concretos para ferramentas de segurança; e (3)~\textit{Parsing Module}, que interpreta as saídas das ferramentas e retroalimenta o raciocínio.
  A avaliação é realizada em desafios da plataforma HackTheBox (HTB) de dificuldades variadas, comparando o PentestGPT com abordagens baseadas em LLMs sem estruturação do fluxo de \textit{pentesting}.

  \textit{Principais Resultados, Contribuições e Direções Futuras:} O PentestGPT supera significativamente a linha de base em quase todas as categorias de desafio, demonstrando que a decomposição estruturada do processo de \textit{pentesting} é compatível com as capacidades atuais dos LLMs.
  Como contribuições, os autores publicam o \textit{framework}, um \textit{benchmark} de avaliação e discutem limitações na execução de tarefas complexas de múltiplos passos.
  As direções futuras indicam integração com ferramentas de segurança especializadas, expansão para cenários de aplicações \textit{web} e adoção de arquiteturas multi-agente para paralelizar subtarefas.

\end{itemize}

\subsection{Trabalho: ReAct: Synergizing Reasoning and Acting in Language Models}

\begin{itemize}

  \item \textbf{Referência:}
  YAO, Shunyu; ZHAO, Jeffrey; YU, Dian; DU, Nan; SHAFRAN, Izhak; NARASIMHAN, Karthik; CAO, Yuan.
  ReAct: Synergizing Reasoning and Acting in Language Models.
  In: \textit{Proceedings of the International Conference on Learning Representations (ICLR)}, 2023.
  Disponível em: \url{https://arxiv.org/abs/2210.03629}.

  \item \textbf{Resumo:}

  \textit{Motivação e Objetivo:} LLMs demonstram capacidade de raciocinar em linguagem natural via \textit{Chain-of-Thought} e de executar ações em ambientes externos via chamadas de ferramenta, mas essas capacidades raramente são exploradas de forma integrada e intercalada.
  O trabalho propõe o padrão ReAct (\textit{Reason + Act}), no qual o modelo alterna, em cada passo, entre a geração de um traço de raciocínio (\textit{thought}) e a execução de uma ação observável, resultando em agentes mais capazes e interpretáveis.

  \textit{Metodologia:} O padrão é avaliado em quatro conjuntos de tarefas: duas de raciocínio baseado em conhecimento (HotpotQA e FEVER) e duas de tomada de decisão interativa em ambientes simulados (ALFWorld e WebShop).
  Para cada tarefa, compara-se ReAct contra \textit{baselines} de raciocínio puro (\textit{Chain-of-Thought}) e de ação pura (\textit{action-only}).
  A implementação utiliza o modelo PaLM (540B parâmetros) com \textit{few-shot prompting}.

  \textit{Principais Resultados, Contribuições e Direções Futuras:} ReAct supera os dois \textit{baselines} em quase todos os \textit{benchmarks}, com ganhos de até 34\% em HotpotQA e melhorias substanciais em tarefas de navegação em ALFWorld e WebShop.
  A principal contribuição é o padrão de \textit{prompt} e o resultado empírico de que intercalar pensamento e ação melhora tanto o desempenho quanto a rastreabilidade do agente.
  Os autores indicam como direções futuras a exploração do padrão em modelos maiores, em modalidades multimodais e combinado com aprendizado por reforço para refinar as estratégias de raciocínio.

\end{itemize}

\subsection{Diferencial}

\textbf{Lacunas identificadas:}

O \textbf{PentestGPT} opera sobre alvos de infraestrutura (servidores, serviços de rede) e plataformas de CTF (\textit{Capture the Flag}), contexto em que a superfície de ataque é relativamente estática e o \textit{feedback} das ferramentas é textual e direto.
O trabalho não aborda o desafio específico das SPAs, cujos estados são gerados dinamicamente por JavaScript e só se tornam acessíveis após interação com o DOM.
Tampouco trata de vulnerabilidades semânticas de lógica de negócio, que exigem conhecimento sobre regras de domínio da aplicação.

O \textbf{ReAct} é uma contribuição de Processamento de Linguagem Natural e não foi concebido para segurança de \textit{software}.
Não propõe nenhuma integração com automação de navegador, protocolos de instrumentação de DOM, ferramentas de análise de segurança ou bibliotecas de casos de abuso.
Trata-se de um padrão metodológico -- valioso como fundação -- mas que, isoladamente, não constitui um \textit{framework} de detecção de vulnerabilidades.

\textbf{Como a proposta se diferencia:}

O \textit{framework} Agentic DAST proposto integra três contribuições que, combinadas, superam as limitações identificadas: (1)~\textbf{navegação semântica de SPAs} via Puppeteer com \textit{Chrome DevTools Protocol} (CDP), tornando visível ao LLM toda a superfície de ataque, incluindo rotas Angular pós-\textit{login}; (2)~\textbf{raciocínio orientado a segurança \textit{web}}, por meio de um \textit{system prompt} especializado que descreve categorias do OWASP Top 10 e orienta o agente a interpretar respostas HTTP como evidências de falhas; e (3)~\textbf{biblioteca de Casos de Abuso}, com oito ataques concretos e determinísticos que garantem cobertura mínima reproduzível em todas as execuções do experimento.
Enquanto o PentestGPT generaliza o \textit{pentesting} para infraestrutura e o ReAct generaliza o padrão de raciocínio para agentes, o Agentic DAST especializa ambas as contribuições para o nicho de SPAs comerciais modernas, avaliando o impacto com métricas formais em um \textit{testbed} padronizado pela comunidade (OWASP Juice Shop).

% ================================================================
\section{Conclusões sobre a Atividade Avaliativa}
% ================================================================

A elaboração desta atividade foi determinante para a maturação da proposta de pesquisa em dois aspectos práticos.

O primeiro diz respeito à \textbf{caracterização do estilo de pesquisa}.
O exercício de classificar o trabalho segundo a taxonomia de Wazlawick (2014) forçou uma reflexão sobre o que, de fato, está sendo produzido: não apenas a implementação de uma ferramenta, mas uma investigação científica com hipóteses mensuráveis e protocolo experimental definido.
Essa distinção é essencial para que o trabalho seja avaliado como contribuição acadêmica -- com afirmações verificáveis -- e não apenas como projeto de engenharia.
Em particular, a definição do procedimento experimental (três cenários, $n = 3$ rodadas, teste de Wilcoxon) surgiu como consequência direta dessa reflexão: sem o enquadramento metodológico, a avaliação poderia ser informal e, portanto, cientificamente fraca.

O segundo aspecto refere-se ao \textbf{posicionamento no estado da arte}.
A leitura crítica dos trabalhos relacionados evidenciou com precisão as lacunas que a proposta deve preencher.
Antes do levantamento, o escopo do trabalho era intuitivo; após identificar que o PentestGPT não trata de SPAs e que o ReAct carece de aplicação em segurança \textit{web}, ficou clara a posição da contribuição na literatura e, consequentemente, quais afirmações o experimento deve sustentar.
Compreender onde o trabalho se situa é condição necessária para formular uma questão de pesquisa respondível e para justificar o esforço experimental diante da comunidade científica.

% ================================================================
\section{Referências}
% ================================================================

\begin{itemize}

  \item DENG, Gelei et al.
  PentestGPT: An LLM-Empowered Automatic Penetration Testing Framework.
  In: \textit{Proceedings of the USENIX Security Symposium}, 2024.
  Disponível em: \url{https://arxiv.org/abs/2308.06782}.

  \item IBM. What is Dynamic Application Security Testing (DAST)?
  \textit{IBM Think}, 2024.
  Disponível em: \url{https://www.ibm.com/think/topics/dast}. Acesso em: mar. 2025.

  \item OWASP FOUNDATION. OWASP Top Ten 2021. 2021.
  Disponível em: \url{https://owasp.org/www-project-top-ten/}. Acesso em: mar. 2025.

  \item PORTSWIGGER. Business Logic Vulnerabilities.
  \textit{Web Security Academy}, 2024.
  Disponível em: \url{https://portswigger.net/web-security/logic-flaws}. Acesso em: mar. 2025.

  \item WAZLAWICK, R. S. \textbf{Metodologia de pesquisa para ciência da computação}. 2. ed.
  Rio de Janeiro: Elsevier, 2014.

  \item YAO, Shunyu et al.
  ReAct: Synergizing Reasoning and Acting in Language Models.
  In: \textit{Proceedings of the International Conference on Learning Representations (ICLR)}, 2023.
  Disponível em: \url{https://arxiv.org/abs/2210.03629}.

\end{itemize}

% Referências internas (apenas para as citações no texto)
\bibliographystyle{plain}
\begin{thebibliography}{9}
\bibitem{yao2023react} YAO et al., ICLR 2023 (ver Seção~2.2).
\bibitem{wazlawick2014} WAZLAWICK, 2014 (ver Seção~4).
\end{thebibliography}

\end{document}
